\section{Information threshold and finite-information precision}\label{information-threshold-and-finite-information-precision}

\textbf{Theorem 2 (Calibration threshold and low information; L.3--L.7).} Uniform consistency on any fixed positive local window is possible if and only if \(M/n\) tends to infinity. If eventually \(M/n\le C_0\) and
\[
0<a\le\min\{A_0,1/(4\sqrt{\max(C_0,1)})\},
\]
then
\[
\liminf_n\mathfrak R_{n,M}(A_0)\ge[2\Phi(ca)-1]/4>0.
\]
If \(M/n\) tends to zero, the exact limiting minimax MAE is \(\Phi(cA_0)-1/2\), attained by the constant forecast .5.

For necessity compare \(\theta=\pm a/\sqrt n\). Product KL is at most \(8a^2M/n\) and total variation at most \(2a\sqrt{M/n}\). Directly integrating absolute loss over the density overlap gives maximum risk at least half the target separation times \(1-\mathrm{TV}\). This avoids the factor of two lost by reducing to a nearest-target testing error. Sufficiency uses \(E_\theta h_0=\theta\) and \(\operatorname{Var}_\theta h_0=1-\theta^2\): the raw moment forecast
\[
\widehat\pi^{\rm mom}=\Phi(c\sqrt n\,\overline h_0)
\]
has uniform error at most \(c\sqrt{n/M}/\sqrt{2\pi}+Ce_n\).

\textbf{Theorem 3 (Finite-ratio minimax limit; L.9--L.10).} If \(M/n\) tends to \(\rho\) in (0,infinity), the entire local calibration experiment on the known tilt subfamily is asymptotically equivalent to \(Y=u+Z/\sqrt\rho\), uniformly for \(|u|\le A_0\). Consequently
\[
\begin{gathered}
\mathfrak R_{n,M}(A_0)\longrightarrow V_G(K,r) \\
K=cA_0 \\
r=c/\sqrt\rho,
\end{gathered}
\]
where
\[
V_G(K,r)=\inf_d\sup_{|t|\le K}E|d(t+rZ)-\Phi(t)|.
\]

The complete experiment-level proof constructs both Markov kernels. Under the baseline, transform the standardized sample moment \(S_M\) by \(T_M=\Phi^{-1}(F_M(S_M))\), with its known continuous baseline CDF \(F_M\). \(T_M\) is exactly standard normal. A uniform log-likelihood expansion and bounded exponential moments give uniform L1 approximation by \(\exp(\sqrt\rho uT_M-\rho u^2/2)\). A baseline conditional-distribution kernel reconstructs the auxiliary full sample from \(T_M\). The two deficiencies tend to zero. Transporting bounded absolute loss and replacing \(\pi_n\) by \(\Phi(cu)\) proves the minimax limit; a scalar CLT alone would not establish it.

For comparison, L.8 proves that the raw moment rule has worst-window limiting risk \(\arctan(r)/\pi\). The constant rule has risk \(\Phi(K)-1/2\). Neither formula equals the general finite-ratio minimax value. Window projection improves the raw rule's loss against \(\Phi(t)\), but does not itself establish optimality.

\subsection{Proof of the information threshold}\label{proof-of-the-information-threshold}

Write \(r_n\) for the uniform error in Theorem 1, so \(r_n\)\textless=Ce\_n.~We retain source equation labels to connect each step to the complete proof chapters.

At \(\theta\)=+/-a/\(\sqrt n\), positivity gives
\[
D_{KL}(f_\theta\|f_{-\theta})\le\chi^2(f_\theta,f_{-\theta})
=4\theta^2 E_0\frac{h_0^2}{1-\theta h_0}\le8\theta^2.
\]
Product KL is at most 8a²\(M/n\). Pinsker gives total variation at most 2a sqrt(\(M/n\))\textless=1/2. Lemma L.2 implies the target separation \(\Delta_n\) tends to Delta=2Phi(ca)-1\textgreater0. Let p\_+,p\_- be the calibration densities and t\_+,t\_- their true exit probabilities. For any forecast T, the triangle inequality on their common density overlap gives
\[
\begin{aligned}
&E_+|T-t_+|+E_-|T-t_-|\\
&\quad\ge\int\min(p_+,p_-)(|T-t_+|+|T-t_-|)\\
&\quad\ge\Delta_n(1-\operatorname{TV}).
\end{aligned}
\]
The maximum risk is at least half this sum, hence at least \(\Delta_n\)/4 when TV\textless=1/2. This proves (L.6). Common risk and fixed source information have identical values and cannot change this comparison. Independent randomization cannot improve the testing bound. \(\square\)

The matching consistency argument uses the raw moment estimator already defined above. Its exact variance is at most 1/M, so the Lipschitz bound gives risk at most c sqrt(n/M)/sqrt(2pi)+Ce\_n.~If the budget ratio does not diverge, a bounded-ratio subsequence and the two-point lower bound exclude consistency.

\subsection{Exact low-information minimax limit}\label{exact-low-information-minimax-limit}

Put \(B_0=\Phi(cA_0)-1/2\). For the upper bound, Theorem 1 and monotonicity of \(\Phi\) give

\[
\sup_{|u|\le A_0}|\tfrac12-\pi_n(Q_{u/\sqrt n})|
\le B_0+r_n.
\]

For the lower bound take a=\(A_0\) in the density-overlap bound. The calibration TV bound tends to zero and the target separation tends to \(2B_0\), so liminf risk\textgreater=\(B_0\). These two bounds agree. If M=0 is allowed, use the identical no-observation experiment, with TV exactly zero. \(\square\)

The endpoint argument establishes the exact asymptotic constant in the regime \(M/n\to0\). The finite-ratio value is instead characterized by Theorem 3. The result strengthens the interpretation of the sample threshold: a best constant predictor may be optimal when information vanishes, while its nonzero worst-case error still prevents uniform consistency.

\subsection{Full-experiment equivalence: likelihoods and both kernels}\label{full-experiment-equivalence-likelihoods-and-both-kernels}

Within this likelihood proof put \(h=h_0\). Write \(P_{n,0}=Q_0^{\otimes M}\), \(\rho_n=M/n\), and

\[
\begin{gathered}
S_n=M^{-1/2}\sum_{j=1}^M h(Z_j) \\
V_n=M^{-1}\sum_{j=1}^M h(Z_j)^2.
\end{gathered}
\]

Let \(F_M\) be the CDF of \(S_n\) under \(P_{n,0}\). In the present family h is a nonconstant real-analytic function, so the preimage of any value is a Lebesgue-null set. Since \(f_0\) has a density, h(Z) has no atoms; the sum also has no atoms. Therefore \(F_M\) is continuous. Define

\[
T_n=\Phi^{-1}(F_M(S_n)).
\]

This parameter-independent statistic has the exact N(0,1) law under \(P_{n,0}\); events where the argument is zero or one have probability zero. The baseline law and the transform are known because the model family is known.

The CLT and continuity of Phi give sup\_x\textbar{}\(F_M\)(x)-Phi(x)\textbar-\textgreater0. On a compact interval for \(S_n\) the inverse normal CDF is uniformly continuous after restricting to a slightly larger compact probability interval. Tightness of \(S_n\) then proves

\[
\begin{gathered}
S_n-T_n\longrightarrow0 \\
\text{in }P_{n,0}\text{-probability}.
\tag{G.2}
\end{gathered}
\]

This use of a CLT is only one part of the proof. The essential next step controls the full likelihood ratios.

Let \(\Lambda_n\)(u)=dP\_\{n,u\}/dP\_\{n,0\}. Boundedness of h and compactness of u allow a uniform Taylor expansion, valid for all sufficiently large n:

\[
\begin{gathered}
\log\Lambda_n(u)
=u\sqrt{\rho_n}S_n-\tfrac12u^2\rho_n V_n+R_n(u) \\
\sup_{|u|\le A_0}|R_n(u)|\le C M n^{-3/2}.
\tag{G.3}
\end{gathered}
\]

Under \(P_{n,0}\), \(V_n\)-\textgreater1 in probability, and \(\rho_n\)-\textgreater{}\(\rho\). Equations (G.2)--(G.3) imply

\[
\sup_{|u|\le A_0}
\left|\log\Lambda_n(u)-\left(\sqrt\rho\,uT_n-\tfrac12\rho u^2\right)\right|
\longrightarrow0
\tag{G.4}
\]

in probability. Define g\_u(t)=exp(sqrt(\(\rho\))u t-\(\rho\) u\^{}2/2). We now strengthen convergence of log likelihoods to uniform L1 convergence of likelihoods.

For any fixed p\textgreater1, \(\log(1+x)\le x\) gives

\[
\sup_{|u|\le A_0}\Lambda_n(u)
\le\exp\left(A_0\left|\sum_jh(Z_j)\right|/\sqrt n\right).
\]

The centered bounded summands have a sub-Gaussian exponential-moment bound. Since \(M/n\) is bounded, the pth moments of the last display are uniformly bounded. Also sup\_u g\_u(\(T_n\))\textless=exp(sqrt(\(\rho\))\(A_0\)\textbar{}\(T_n\)\textbar), whose pth moment is bounded because \(T_n\) is exactly standard normal. Equation (G.4), tightness, and these bounds give uniform integrability and hence

\[
\varepsilon_n:=\tfrac12\mathbb E_{n,0}
\sup_{|u|\le A_0}|\Lambda_n(u)-g_u(T_n)|\longrightarrow0.
\tag{G.5}
\]

Define an auxiliary experiment on the original sample space by

\[
d\widetilde P_{n,u}=g_u(T_n)\,dP_{n,0}.
\]

These are probability measures, because E\_\{n,0\}g\_u(\(T_n\))=1 exactly. Equation (G.5) bounds sup\_u TV(\(P_{n,u}\),\(\widetilde P_{n,u}\)) by \(\varepsilon_n\).

Under \(\widetilde P_{n,u}\), \(T_n\) has law N(sqrt(\(\rho\))u,1). Conversely, use the regular conditional distribution \(P_{n,0}\)(d z\textbar{}\(T_n\)=t) to turn such a normal observation back into a full calibration sample. This kernel is independent of u; tilting by a function of \(T_n\) leaves the conditional law unchanged. Regular conditional distributions exist on the Euclidean sample spaces, and arbitrary versions on baseline-null sets do not matter because all normal shifts are mutually absolutely continuous. Thus the auxiliary experiment is exactly equivalent to observing T\textasciitilde N(sqrt(\(\rho\))u,1), or Y=T/\(\sqrt\rho\).

Applying the statistic \(T_n\) in one direction and this baseline conditional kernel in the other gives both deficiencies bounded by \(\varepsilon_n\). This proves the full-experiment equivalence in Theorem 3. \(\square\)

The proof supplies asymptotic equivalence on a fixed compact local parameter interval. It does not give a calibrated numerical value of \(\varepsilon_n\), nor global equivalence over an expanding parameter range. The conditional kernel is an existence construction; implementing it is not required to use a decision rule in the forward direction.

\subsection{Minimax risk transfer}\label{minimax-risk-transfer}

Project all decision values to {[}0,1{]}, which cannot increase loss. Absolute loss is then bounded by one, so transporting any randomized rule through either Markov kernel changes each parameterwise risk by at most \(\varepsilon_n\). Replacing the target \(\pi_n\) by \(\Phi(cu)\) changes each loss by at most \(r_n\). Taking supremum risks and then infima in both directions gives

\[
|\mathfrak R_{n,M}(A_0)-V(A_0,c,\rho)|\le\varepsilon_n+r_n.
\]

This reasoning does not require the minimax infimum to be attained: apply it to arbitrarily near-optimal rules. Under absolute loss, averaging any external decision randomization conditionally on the observation cannot worsen risk, by convexity. Thus restricting to deterministic rules leaves the values unchanged. \(\square\)

Standardize with

\[
\begin{gathered}
t=cu \\
K=cA_0 \\
X=cY=t+rZ \\
r=c/\sqrt\rho.
\end{gathered}
\]

The limiting decision problem depends only on (K,r):

\[
V_G(K,r)=\inf_d\sup_{|t|\le K}\mathbb E_t|d(X)-\Phi(t)|.
\tag{G.7}
\]

Actions may be restricted further to {[}\(\Phi(-K)\),\(\Phi(K)\){]}. Reflection and averaging permit a symmetric rule d(-x)=1-d(x) without increasing worst-case risk, although the certificates below do not require a proof that the candidate optimizer is minimax.

\subsection{Uniform risk of the raw moment predictor}\label{uniform-risk-of-the-raw-moment-predictor}

By the exact calibration moment identities, the standardized calibration sum

\[
Z_{n,u}=\frac{\sqrt M(\widehat\theta-u/\sqrt n)}
 {\sqrt{1-u^2/n}}
\]

has a uniform Berry--Esseen bound C/\(\sqrt M\) on the compact u-window: h is bounded and its variance is uniformly bounded away from zero for large n.~Hence

\[
c\sqrt n\widehat\theta
=cu+c\sqrt{n/M}\sqrt{1-u^2/n}\,Z_{n,u}.
\]

The coefficient of \(Z_{n,u}\) tends uniformly to r. This implies uniform convergence of the bounded loss expectations to \(D(cu,r)\). One direct justification is by subsequences: any alleged nonuniformity supplies u\_n in the compact window; extract u\_n-\textgreater u\_*, apply the triangular CLT and bounded continuous mapping along that subsequence, and use continuity of D. Replacing \(\Phi(cu)\) by the true \(\pi_n\) changes expected absolute loss by at most \(r_n\). This proves the uniform moment-risk convergence.

Pair positive and negative Z to obtain

\[
D(t,r)=\int_0^\infty
 [\Phi(t+rz)-\Phi(t-rz)]\varphi(z)\,dz.
\tag{L.18}
\]

Symmetry makes this even in t. For t\textgreater0 and z\textgreater0,
\textbar t+rz\textbar\textgreater\textbar t-rz\textbar, so differentiation under the integral gives a strictly negative derivative in t. Therefore the maximum over a window containing zero is \(D(0,r)\).

Finally,

\[
\frac{\partial}{\partial r}D(0,r)
=2\int_0^\infty z\varphi(rz)\varphi(z)\,dz
=\frac1{\pi(1+r^2)}.
\]

Since D(0,0)=0, integration yields \(D(0,r)\)=arctan(r)/pi. Uniform convergence in the uniform moment-risk convergence allows passage to the supremum, proving (L.17). \(\square\)

The arctan expression remains a result for this specified predictor. It is now its worst-case local-window risk limit, rather than a universal lower bound. In this known family, the constant predictor has worst-case limiting risk \(B_0\); it can outperform the raw moment predictor for some combinations of window width and calibration ratio.
